\chapter{Permanent Pacemaker Implantation}
\section*{Overview}
A permanent pacemaker is an implanted device that provides electrical impulses when the heart beats too slowly or irregularly.
\section*{Why It Is Done}
It is used for symptomatic bradycardia, heart block, sick sinus syndrome, and selected rhythm disorders.
\section*{How It Is Performed}
A generator is placed beneath the skin below the collarbone and one or more leads are guided through a vein into the heart. The system is tested and programmed before closure.
\section*{Risks and Follow-up}
Infection, bleeding, pneumothorax, lead displacement, vessel or heart injury, and device malfunction are possible. Regular device checks and eventual generator replacement are required.
\section*{References}
\begin{enumerate}\item MedlinePlus. Pacemaker. U.S. National Library of Medicine; 2025.\item Current Medical Diagnosis and Treatment. McGraw-Hill; 2025.\end{enumerate}
\clearpage
